\section*{Deviations from native Python}

We intend the Python \S $x$ to be a conservative extension of 
native Python: Every correct Python \S $x$ program should behave 
\emph{exactly} the same using a Python \S $x$ implementation, as it does 
using a native Python implementation. We assume, of course, that 
suitable libraries are used by the TypeScript implementation, to 
account for the predefined names of Python \S $x$. 
This section lists some exceptions where we think a Python \S $x$ 
implementation should be allowed to deviate from the native Python 
specification, for the sake of internal consistency and esthetics.

\begin{description}  
\item[{Output Behavior Differences Between py-slang and Standard Python REPL:}]
  Python provides an interactive mode, where users can type and immediately 
  evaluate Python expressions line by line. This mode is entered when you 
  simply type python in a terminal without specifying a file. This is commonly 
  referred to as the Python REPL. 
  In the standard Python REPL, any evaluated expression automatically has its 
  result printed to the console, even if the user does not explicitly call 
  \texttt{print()}. For example:

  \begin{lstlisting}
>>> 1 + 2
3
>>> "hello"
'hello'
  \end{lstlisting}

  In contrast, py-slang implements the behavior of \textbf{non-interactive}
  Python, the mode used when you pass a Python program to \texttt{python}.
  Only expressions explicitly passed to \texttt{print()} produce output. If the 
  expression is evaluated without a print call, no output will appear, even 
  though a value is computed.

  For example, in py-slang:

  \begin{lstlisting}
print(1 + 2)  	# Outputs: 3
  \end{lstlisting}

  But

  \begin{lstlisting}
1 + 2    	# Outputs nothing
  \end{lstlisting}

\item[{Tail calls:}]
  A call in tail position runs in constant space, so that a recursive function can
  run for any number of iterations: \texttt{f(10000)} with a tail-recursive
  \texttt{f} works, where native Python raises a \texttt{RecursionError} at a depth
  of about 1000.

\item[{Conditions and logical operators:}]
  The conditions of \texttt{if}, \texttt{elif} and \texttt{while}, and the operands
  of \texttt{not}, and the first operand of \texttt{and} and \texttt{or}, must be
  booleans, as the tables of operator types specify; native Python accepts any value
  (truthiness). \texttt{if 1:} and \texttt{0 or 'x'} are errors.

\item[{Dynamic type checks:}]
  The operators only accept the types of the tables of operator types. Native
  Python accepts more: \texttt{True + True}, \texttt{'ab' * 3}, \texttt{[1] < [2]}
  and \texttt{min([3, 1, 2])} are errors here.

\item[{Numbers inside lists:}]
  Numbers are compared by value, also when they are inside a list or pair, so that a
  \texttt{nan} in a list equals every \texttt{nan} in the other list, also one that
  was computed separately (native Python says False for two nans that were computed
  separately).

\item[{Limits of some implementations:}]
  These differences from native Python are due to the current state of the implementations,
  not to the language, and will be removed.
  \begin{itemize}
  \item A power with a float exponent that is not an integer can differ from Python's
    in the last digit, and a negative base with a fractional exponent gives \texttt{nan}
    instead of a complex number.
  \item The escape sequence \texttt{\textbackslash U} of strings is not read.
  \item The implementations that run on the Pynter virtual machine (PVML, native Pynter)
    and the WebAssembly implementation evaluate the list, the index and then the value
    of \texttt{list[index] = value}; native Python evaluates the value first.
  \item The WebAssembly implementation has 64-bit integers, a lower limit for the
    depth of recursion, and does not support \texttt{global} in nested functions.
  \item The environment stepper (\S3 and \S4) raises a \texttt{RecursionError} when 300 calls
    are waiting for their result, where native Python allows about 1000.
  \end{itemize}
\end{description}
